\documentclass{article}
\title{Circumscribe}
\author{Rudomotkin Ivan}
\date{November 14, 2026}
\usepackage{amsmath}
\renewcommand{\arraystretch}{1.5}
\begin{document}

\maketitle

\tableofcontents

\section{Introduction}

`Circumscribe' is a general-purpose, language-agnostic macro templating tool.
Its purpose is recursively finding and subsituting all placeholders within the
provided file or directory.

\subsection{Substitution}

The \texttt{circumscribe substitute} subcommand is responsible for all substitution.

\section{Language}

Circumscribe features a domain-specific, expression-first, \emph{statically typed}, functional programming language.
Unlike most general-purpose languages, circumscribe does not have a notion of a
`type': \emph{The only type in circumscribe is a string}. Circumscribe was designed
with this constraint in mind.

In order to compensate for the lack of `types', circumscribe offers a notion of \emph{patterns},
a way to denote a specific formal language. Patterns completely replace the type system of circumscribe.

A circumscribe source file may contain single-line comments beginning with a hash sign. A comment may appear anywhere in
the source file. The language is whitespace-insensitive. It's a convention to end a circumscribe source file in \texttt{.circumscribe}.

\subsection{Snippet}

A \emph{snippet} is a piece of text. It is analogous to a string, and both terms are often used interchangeably.
A snippet is the only type in circumscribe.

\subsection{Macro}

A \emph{macro} is a single reusable piece of code that inevitably produces a \emph{snippet}. It is analogous to a function,
and both terms are often used interchangeably. A macro is defined with the \texttt{def} keyword at the topmost level of
a circumscribe source file.

\begin{center}
    \verb|def <name>(<name>: <pattern>, ...) -> <return>: <expression>|
\end{center}

A macro may take in an arbitrary number of parameters in the form of \verb|<name>: <pattern>|.
The pattern is optional and may be omitted, in which case it defaults to \texttt{any}.
The arguments passed to the macro may be referenced by their name.


\begin{center}
    \verb/def greet(name): |Hello, <name>!|/
\end{center}

Circumscribe does not differentiate between macros and ``variables''.
A ``variable'' is merely a nullary macro. However, the parameter list
is itself optional and may too be omitted:

\begin{center}
    \verb/def greeting: greet(|Marie|)/
\end{center}

A macro may have its return pattern annotated. This, however, isn't mandatory.
As long as the macro doesn't recurse, circumscribe is capable of inferring the
return pattern manually, as observed in above examples.

Circumscribe differentiates between pure and impure macros. A macro is considered
impure if its body refers to an impure macro, and pure until proven otherwise.

\subsection{Expression}

An expression is a single piece of code that may be evaluated to a \emph{snippet}.
Every expression \textbf{must} evaluate to a snippet. Because of this promise, all
control-flow happens in terms of expressions.

\subsubsection{Literal}

A \texttt{literal} represents a piece of literal text, enclosed in a pair of vertical bars.
Everything in-between the delimiters is taken literally from the source.
A literal may span across multiple lines.

\begin{center}
    \verb/|Hello, Marie!|/
\end{center}

A literal may include escape sequences. An escape sequence is denoted by a backslash, followed
by an escape character.

\begin{center}
\begin{tabular}{ c|l r }
    \hline
    Escape Sequence & Replacement & Unicode \\[1ex]
    \hline
    \verb/\|/ & Vertical Bar    & \texttt{u007C} \\ 
    \verb|\<| & Less Than       & \texttt{u003C} \\
    \verb|\\| & Backslash       & \texttt{u005C} \\
    \verb|\n| & New Line        & \texttt{u000A} \\ 
    \verb|\r| & Carriage Return & \texttt{u000D} \\ 
    \verb|\t| & Tabulation      & \texttt{u0009} \\ 
    \verb|\b| & Backspace       & \texttt{u0008} \\
\end{tabular}
\end{center}

Alternatively, an escape sequence may be of form \verb|\uXXXX|, in which case it is substituted for
the literal Unicode point.

\begin{center}
    \verb/|Hello, \u004D\u0061\u0072\u0069\u0065!|/
\end{center}

A literal may contain placeholders, in the form of \texttt{<(expression)>}. When the literal is evaluated, all
placeholders are subtituted and concatenated to the rest of the literal. This is commonly referred to as
\emph{snippet interpolation}.

\begin{center}
    \verb/|Hello, <name>!|/
\end{center}

The inferred pattern of a \texttt{literal} is the product pattern of all constituents of the literal.

\subsubsection{Application}

An \texttt{application} represents a single application of a macro.

\begin{center}
    \texttt{<name>($arg_0$, $arg_1$, \dots, $arg_n$)}
\end{center}

In order to compensate for the lack of variables, circumscribe has a notion
of a \emph{bare} application. An application is considered bare if the argument
list is omitted. A \emph{bare} application is only possible on a \emph{pure},
\emph{nullary} macro.

\begin{center}
    \verb/|Hello, <name>!|/
\end{center}

The inferred pattern of an \texttt{application} is the return pattern of the macro.

\subsubsection{If}

An \texttt{if} decides on a branch based on a condition. The condition must satisfy the standard
\texttt{boolean} pattern. Since all expressions must evaluate to a snippet, an else branch
is mandatory.

\begin{center}
    \verb|if <condition>: <then> else <else>|
\end{center}

The inferred pattern of an \texttt{if} is the sum pattern of both branches.

\subsubsection{Match}

A \texttt{match} matches a source against a \emph{grammar}, then evaluates
the corresponding branch. If the source could successfully match the grammar,
the named placeholders may be referenced in the body of the expression. Since all expressions must evaluate to a snippet, an else branch
is mandatory.

\begin{center}
    \verb|match <grammar> in <source>: <then> else <else>|
\end{center}

The inferred pattern of a \texttt{match} is the sum pattern of both branches.

\subsubsection{Compare}

A \texttt{compare} determines whether the snippet produced by the left-hand side of the expressions
is identical to the one produced by the right-hand side of the expression.

\begin{center}
    \verb|<left> = <right>|
\end{center}

Due to circumscribe being a functional language and having no notion of an `assignment',
we may use the \texttt{=} symbol directly. The inferred pattern of a \texttt{compare} is the
standard \texttt{boolean} pattern.

\subsubsection{And}

An \texttt{and} represents the logical conjunction on two expressions. Both expressions must
satisfy the standard \texttt{boolean} pattern. If the left-hand side of the expression
evaluates to \texttt{false}, the whole expression short-circuits immediately.

\begin{center}
    \verb|<left> and <right>|
\end{center}

The inferred pattern of an \texttt{and} is the standard \texttt{boolean} pattern.

\subsubsection{Or}

An \texttt{or} represents the logical disjunction on two expressions. Both expressions must
satisfy the standard \texttt{boolean} pattern. If the left-hand side of the expression
evaluates to \texttt{true}, the whole expression short-circuits immediately.

\begin{center}
    \verb|<left> or <right>|
\end{center}

The inferred pattern of an \texttt{or} is the standard \texttt{boolean} pattern.

\subsubsection{Not}

A \texttt{not} represents the logical negation on an expression. The expression must satisfy 
the standard \texttt{boolean} pattern.

\begin{center}
    \verb|not <expression>|
\end{center}

The inferred pattern of a \texttt{not} is the standard \texttt{boolean} pattern.

\subsection{Pattern}

A \texttt{pattern} represents a single formal language. It replaces the notion of a
`type' and serves an identical duty. In order for a pattern to accept a different one,
the latter has to be a subset of the former. We say the former \emph{satisfies}
the latter.

\[L_0\left(a\right) \subseteq L_1\left(b\right) \iff a\text{ satisfies }b\]

A pattern may be aliased with the \texttt{pat} keyword at the topmost level of a
circumscribe source file.

\begin{center}
    \verb/pat <name>: <pattern>/
\end{center}

A pattern may not recurse. A recursive formal language is no longer regular, and
no meaningful operations may be formed on context-free languages.

\begin{center}
    \verb/pat balanced: |(| * (balanced + boolean) |)|/
\end{center}

We say a pattern is \emph{nullable} if it accepts an empty snippet.

\subsubsection{Literal}

A literal accepts only a single literal snippet. A literal is enclosed in a pair
of vertical bars, visually mirroring the syntax of the literal expression.

\begin{center}
    \verb/|Hello, Marie!|/
\end{center}

A literal may contain escape sequences, similarly to the literal expression.

\begin{center}
\begin{tabular}{ c|l r }
    \hline
    Escape Sequence & Replacement & Unicode \\[1ex]
    \hline
    \verb/\|/ & Vertical Bar    & \texttt{u007C} \\ 
    \verb|\\| & Backslash       & \texttt{u005C} \\
    \verb|\n| & New Line        & \texttt{u000A} \\ 
    \verb|\r| & Carriage Return & \texttt{u000D} \\ 
    \verb|\t| & Tabulation      & \texttt{u0009} \\ 
    \verb|\b| & Backspace       & \texttt{u0008} \\
\end{tabular}
\end{center}

It may, however, not contain placeholders. Instead, patterns may be concatenated
with multiplication:

\begin{center}
    \verb/|Hello, | * any * |!|/
\end{center}

\subsubsection{Interval}

An interval accepts a single character anywhere in-between its bounds, inclusively.
Two literals may be converted into an interval with a colon, as long as each one
satisfies the \verb/|\u0000|-|\u10FFFF|/, the \verb|w| pattern.

\begin{center}
    \verb/|0|:|9|/
\end{center}

\subsubsection{Sum}

A \texttt{sum} represents a disjunction of two formal languages. Sum is analogous to
an alternation in the usual regular expression sense.

\[a + b = L\left(a\right) \cup L\left(b\right)\]

\subsubsection{Product}

A \texttt{product} represents a cartesian product of two formal languages. Product is analogous
to an juxtaposition in the usual regular expression sense.

\[a * b = L\left(a\right) \times L\left(b\right)\]

\subsubsection{Ellipsis}

An \texttt{ellipsis} represents an arbitrary repetition of a formal language.
Ellipsis is analogous to a Kleene's star in the usual regular expression sense.
The subpattern may not be nullable.

\[a... = L\left(a\right)^{*}\]

\subsubsection{Intersection}

An \texttt{intersection} represents the intersection of two formal languages.
Intersection has no counterpart in the usual regular expression sense.

\[a \text{ \& } b = L\left(a\right) \cap L\left(b\right)\]

\subsubsection{Complement}

A \texttt{complement} represents the negation of a formal languages.
Complement has no counterpart in the usual regular expression sense.

\[!a = \neg L\left(a\right)\]

\subsubsection{Reference}

A \texttt{reference} refers to a pattern alias, and is identical to the aliased pattern.


\begin{center}
    \verb/def foo() -> boolean: |true|/
\end{center}

\subsection{Grammar}

A grammar describes the structure of a snippet. A grammar visually resembles a literal expression,
and may include named placeholders. A grammar may be \emph{matched} against a snippet, capturing
the content of the placeholders.

A grammar is more commonly known as a \emph{pattern}, and the process of matching a grammar
as \emph{pattern matching}. However, circumscribe differentiates between the two terms.

\begin{center}
    \verb/|Hello, <name: any>!|/
\end{center}

A placeholder needn't be named. It is, however, often desired to capture a part of the snippet. The snippet is matched
against the grammar greedily, with backtracking.

\begin{center}
    \verb/|Hello,<s...><name: any>!|/
\end{center}

A grammar is identical to the product pattern of all its literal constituents
and placeholders. Any pattern is a grammar, though a grammar isn't a pattern. In case your grammar consists of a
single pattern, you may write the pattern directly:

\begin{center}
    \[\verb/|<boolean>|/ \cong \verb/boolean/\] 
\end{center}

A grammar may contain escape sequences, similarly to the literal expression.

\begin{center}
\begin{tabular}{ c|l r }
    \hline
    Escape Sequence & Replacement & Unicode \\[1ex]
    \hline
    \verb/\|/ & Vertical Bar    & \texttt{u007C} \\ 
    \verb|\<| & Less Than       & \texttt{u003C} \\
    \verb|\\| & Backslash       & \texttt{u005C} \\
    \verb|\n| & New Line        & \texttt{u000A} \\ 
    \verb|\r| & Carriage Return & \texttt{u000D} \\ 
    \verb|\t| & Tabulation      & \texttt{u0009} \\ 
    \verb|\b| & Backspace       & \texttt{u0008} \\
\end{tabular}
\end{center}

\end{document}